\section{Introduction}\label{intro}
Often, past information carried on by time series data can help one understand future phenomena. Time series forecasting is hereafter attractive in many research fields and lots of research work are developed to improve the prediction ability of the models. Traditionally, autoregressive models and moving average models are often applied to time series prediction. However, they essentially only capture linear relationship in the data set and can not afford many real cases, such as weather forecasting, financial market prediction or signal processing, which are often complex, nonlinear or even chaotic. Another class of statistical methods is Bayesian inference which is leveraged to deal with uncertainty in noisy environment and improve model accuracy by collecting some prior information. For example, Sahoo and Patra \cite{Sahoo18092020} combine the Bayesian inference method with Markov Chain Monte Carlo (MCMC) algorithm to evaluate water pollution rates. With the help of some prior knowledge, this method can greatly improve the prediction accuracy. However, if the prior distribution is not presumed correctly, the inference will go to the wrong direction. Hence, improving the reliability of prior knowledge is crucial to the model performance. As
there always exists various kinds of uncertainties in the forecasting problems, how to accurately capture the intrinsic volatile nature of time series data is still challenging.

Recently, machine learning models outperform conventional methods in many aspects and have led a growing amount of researches to utilize various deep learning frameworks in the research of time series forecasting. Nowadays, the partnership between deep learning and dynamical systems has become a topic of increasingly popularity, as they could benefiting each other in both theoretical and computational aspects. On the computational side, for example, Neural ODE \cite{chen2019neuralordinarydifferentialequations} is a combination of residual network and ordinary differential equations. For stochastic cases, neural stochastic differential equation (Neural SDE) \cite{liu2019neuralsdestabilizingneural} shows a continuous neural network framework based on stochastic differential equation with Brownian motions.  The core of the SDE-Net is to treat the deep neural network transformation as the state evolution of a stochastic dynamic system, and introduce the Brownian motion term to capture cognitive uncertainty.

In many real world applications, Lévy noise plays an important part as a more general random fluctuation including complex non-Gaussian properties. Zhu \cite{ZHU2014126} compares the Lévy noise with the standard Gaussian noise and concludes that the Lévy noise is more versatile and has a wider range of applications in diverse areas. The properties of Lévy motion enables it to better capture the uncertainty in financial data, such as abrupt transitions.

\newpage

The rest of this report is organized as follows. In Section 2, Neural network models based on stochastic differential networks induced by Brownian motion and Lévy motions are presented. In Section 3, experimental results with real financial data are reported. Finally, I conclude the report and give future scopes of work in Section 4.
